\documentclass[10pt,a4paper]{article}

\usepackage[utf8]{inputenc}
\usepackage[T1]{fontenc}
\usepackage[english]{babel}
\usepackage{amsmath,amssymb}
\usepackage[top=1.1cm,bottom=1.1cm,left=1.7cm,right=1.7cm]{geometry}
\usepackage{titlesec}
\usepackage{hyperref}
\hypersetup{colorlinks=true,linkcolor=black,citecolor=blue,urlcolor=blue}

\setlength{\parskip}{0.15em}
\titlespacing*{\section}{0pt}{0.7em}{0.3em}

\title{\vspace{-1.2cm}Bibliographic Review, Topic 39\\[0.2em]
\large The Blind Spot Paradox}
\author{Alexandre Boyer \and Mélaine Gouillou \and Salomé Fonvielle \and Ulysse Petit-Tichanné\\[0.2em]
\normalsize Research Track, CentraleSupélec. Supervisor: Raphaël Minato}
\date{}

\begin{document}
\maketitle
\vspace{-1.4em}

\section{Defining and classifying drift}

The vocabulary was unstable for years, and Moreno-Torres et al.~\cite{moreno2012} write
their taxonomy to fix it: covariate shift when only $P(X)$ moves, prior probability shift
when $P(y)$ does, and concept shift, the case usually called concept drift, when
$P(y\mid X)$ changes. The contribution is terminological rather than algorithmic, and they
say so.

Two surveys map the field. Gama et al.~\cite{gama2014} organise adaptation into forgetting
mechanisms, ensembles and explicit detectors; Lu et al.~\cite{lu2018} restructure the same
material around detection, understanding and adaptation, and classify drift by shape.
Both work from one assumption: the detector serves the classifier, the classifier reacts to
the detector, and neither framework anticipates the two working against each other.

\section{Adaptive learners on streams}

Domingos and Hulten~\cite{domingos2000} introduced VFDT, a decision tree that reads each
example once in constant time and memory, using a Hoeffding bound to decide when a node has
seen enough data to split with confidence. It agrees asymptotically with a batch learner on
stationary data, and it never undoes a split, so it does not adapt.

ADWIN~\cite{bifet2007} adds the adaptivity VFDT lacks. It keeps a window of variable
length, cuts it at every
possible point, and declares a change when two sub-window means differ by more than a
Hoeffding bound. Unlike most competing heuristics it comes with proven bounds on
false positive and false negative rates, and its second version runs in $O(\log W)$ memory
and update time.

The Adaptive Random Forest~\cite{gomes2017} combines the two: $M$ Hoeffding trees,
Poisson$(\lambda{=}6)$ resampling as in leveraging bagging, and per-tree detection with two
thresholds, a permissive $\delta_w$ that starts a background tree and a stricter $\delta_d$
that promotes it. ARF is evaluated on predictive accuracy alone, as is conventional for
ensembles of this kind. River~\cite{montiel2021} is the Python library implementing these
algorithms, and its defaults are what deployed pipelines inherit.

\section{Detecting change from an error stream}

Page~\cite{page1954} proposed the cumulative sum for industrial quality control, against
moving averages whose window length forces a choice between reacting to large shifts and
detecting small ones. CUSUM accumulates deviations from a reference level and fires when the
running sum passes a threshold, so its delay grows as that threshold divided by the size of
the shift. Hinkley~\cite{hinkley1971} takes the same statistic and asks where the change
occurred rather than whether it did, deriving the asymptotic law of that estimate for normal
observations with known initial mean and variance. Both assume a change that persists.

KSWIN~\cite{raab2020} belongs to another family, comparing two sliding windows with a
non-parametric Kolmogorov-Smirnov test and pairing detection with prototype-based
adaptation. Its authors report more false positives than ADWIN or DDM, and argue they are
harmless because the adaptation they trigger is cheap.

\section{Evaluation and synthetic ground truth}

Gama et al.~\cite{gama2014} describe three protocols: holdout, which needs a test set drawn
from the concept active at time $t$ and is therefore restricted to synthetic data;
prequential, where each instance is tested before training, over a landmark window, a
sliding window or fading factors; and controlled permutations, which rerun the stream in
randomised but locality-preserving orders. Lu et al.~\cite{lu2018} add the metrics: kappa
and kappa-temporal for imbalanced or temporally dependent labels, prequential AUC,
RAM-hours for cost. Neither lists a measure of how long a model takes to recover.

Ground truth is the obstacle in finance, since nobody knows when a real regime actually
changed. ProteuS~\cite{suarez2025} fits ARMA-GARCH models to four ETFs and simulates abrupt
and gradual transitions between the fitted regimes with known break dates. Its authors' own
analysis finds heavy overlap between the market states, so the task stays hard.

\section{The source paper}

\emph{The Blind Spot Paradox}~\cite{blindspot2026} attacks the cooperative assumption.
Instrumenting the ARF's swap counter in River 0.23.0, it argues that an adaptive classifier
can erase the evidence an external detector needs. The first mechanism is a clock artefact:
River evaluates ADWIN every $c$ steps, 32 by default, while the ARF's internal detectors run
at $c=1$, which handicaps the external monitor by construction. The second, named the Hydra
effect, is that the forest adapts as soon as its fastest tree does, measured 4 to 8 times
faster than a single tree. The third is starvation: the transient error excess and the CUSUM
accumulation time both scale as $O(1/\Delta e)$, so a violent shock is repaired fast enough
to leave no more evidence than a slow one. On heteroscedastic ProteuS streams a Page-Hinkley
test paired with an ARF at $c_{\mathrm{int}}{=}1$ detects none of 1{,}080 runs. The proposed
Decoupling Principle asks that the external monitor be strictly more reactive than the
internal one, and the paper reports that on those streams this conflicts with false-alarm
control, leaving no admissible threshold. A windowed KSWIN monitor escapes the effect
altogether.

The paper is explicit about what it leaves unsettled. The starvation boundary assumes
conditional independence between trees, which the shared stream violates, and the authors
note this is why the measured acceleration falls short of the $M$-fold reference. Swap
tracking mixes genuine detections with background noise, which rules out isolated
$\mathrm{ARL}_0$ comparisons. And the predicted starvation appears only on stationary
synthetic streams: on real non-stationary data the failure inverts into a collapse of
precision through false-alarm flooding.

\section{Statistical background}

For the law of $\min_i \tau_i$ the paper draws on classical order statistics~\cite{david2003},
whose standard result gives the expected minimum of $M$ independent variables as
$\int_0^\infty (1-F)^M$. That says nothing once the variables are dependent. Esary, Proschan
and Walkup~\cite{esary1967} handle that case, defining association through
$\mathrm{Cov}[f(T),g(T)]\ge 0$ for all non-decreasing $f$ and $g$, and proving that
associated variables satisfy $\Pr(\tau_1>s_1,\dots,\tau_M>s_M)\ge\prod_i\Pr(\tau_i>s_i)$,
the reverse of the independent case.

A separate literature deals with two competing stopping times. Fine and Gray~\cite{fine1999}
model the cumulative incidence function $F_1(t)=\Pr(T\le t,\ \varepsilon=1)$, the probability
of failing from a given cause by $t$, arguing that the cause-specific hazard has no direct
reading in those terms; their own caveat is that the corresponding risk set is artificial,
since it keeps subjects who have already failed from another cause. Aalen and
Johansen~\cite{aalen1978} give the estimator this needs, a product-limit estimator of the
transition probabilities of a non-homogeneous Markov chain under censoring.

\vspace{-0.2em}
{\footnotesize
\begin{thebibliography}{99}\setlength{\itemsep}{0pt}
\bibitem{moreno2012} J.~G. Moreno-Torres, T.~Raeder, R.~Alaiz-Rodríguez, N.~V. Chawla, and F.~Herrera, ``A unifying view on dataset shift in classification,'' \emph{Pattern Recognition}, vol.~45, no.~1, pp.~521--530, 2012.
\bibitem{gama2014} J.~Gama, I.~Žliobaitė, A.~Bifet, M.~Pechenizkiy, and A.~Bouchachia, ``A survey on concept drift adaptation,'' \emph{ACM Computing Surveys}, vol.~46, no.~4, art.~44, 2014.
\bibitem{lu2018} J.~Lu, A.~Liu, F.~Dong, F.~Gu, J.~Gama, and G.~Zhang, ``Learning under concept drift: a review,'' \emph{IEEE Trans. Knowl. Data Eng.}, vol.~31, no.~12, pp.~2346--2363, 2018.
\bibitem{domingos2000} P.~Domingos and G.~Hulten, ``Mining high-speed data streams,'' in \emph{Proc. KDD}, 2000, pp.~71--80.
\bibitem{bifet2007} A.~Bifet and R.~Gavaldà, ``Learning from time-changing data with adaptive windowing,'' in \emph{Proc. SIAM SDM}, 2007, pp.~443--448.
\bibitem{gomes2017} H.~M. Gomes, A.~Bifet, J.~Read, J.~P. Barddal, F.~Enembreck, B.~Pfahringer, G.~Holmes, and T.~Abdessalem, ``Adaptive random forests for evolving data stream classification,'' \emph{Machine Learning}, vol.~106, no.~9--10, pp.~1469--1495, 2017.
\bibitem{montiel2021} J.~Montiel, M.~Halford, S.~M. Mastelini et al., ``River: machine learning for streaming data in Python,'' \emph{JMLR}, vol.~22, pp.~1--8, 2021.
\bibitem{page1954} E.~S. Page, ``Continuous inspection schemes,'' \emph{Biometrika}, vol.~41, no.~1--2, pp.~100--115, 1954.
\bibitem{hinkley1971} D.~V. Hinkley, ``Inference about the change-point from cumulative sum tests,'' \emph{Biometrika}, vol.~58, no.~3, pp.~509--523, 1971.
\bibitem{raab2020} C.~Raab, M.~Heusinger, and F.-M. Schleif, ``Reactive soft prototype computing for concept drift streams,'' \emph{Neurocomputing}, vol.~416, pp.~340--351, 2020.
\bibitem{suarez2025} A.~L. Suárez-Cetrulo, A.~Cervantes, and D.~Quintana, ``ProteuS: a generative approach for simulating concept drift in financial markets,'' arXiv:2509.11844, 2025.
\bibitem{blindspot2026} Anonymous, ``The blind spot paradox: when adaptive classifiers defeat drift detectors,'' under triple-blind review, IEEE ICDM 2026.
\bibitem{david2003} H.~A. David and H.~N. Nagaraja, \emph{Order Statistics}, 3rd ed. Wiley, 2003. (Not available online; the association results of~\cite{esary1967} were consulted instead.)
\bibitem{esary1967} J.~D. Esary, F.~Proschan, and D.~J. Walkup, ``Association of random variables, with applications,'' \emph{Ann. Math. Statist.}, vol.~38, no.~5, pp.~1466--1474, 1967.
\bibitem{fine1999} J.~P. Fine and R.~J. Gray, ``A proportional hazards model for the subdistribution of a competing risk,'' \emph{JASA}, vol.~94, no.~446, pp.~496--509, 1999.
\bibitem{aalen1978} O.~O. Aalen and S.~Johansen, ``An empirical transition matrix for non-homogeneous Markov chains based on censored observations,'' \emph{Scand. J. Statist.}, vol.~5, no.~3, pp.~141--150, 1978.
\end{thebibliography}
}

\end{document}
